\chapter{Нормативно-правове та стандартизаційне забезпечення протидії злоякісному коду}
\chaptershort{Нормативно-правове та стандартизаційне забезпечення}
\label{chap:ch02}

Технічна протидія злоякісному коду діє в межах правового й нормативного поля, що визначає, які дії є протиправними, хто й за що відповідає, які вимоги обов'язкові для організацій і як узгоджуються засоби захисту. Для фахівця важливо розуміти не лише окремі документи, а й те, як вони взаємодіють: кримінальне право встановлює відповідальність за створення й поширення ЗПК; законодавство про кібербезпеку й критичну інфраструктуру зобов'язує організації впроваджувати захист і повідомляти про інциденти; стандарти управління конкретизують процеси; технічні бази знань забезпечують спільну мову опису загроз і контрзаходів. Розділ систематизує правове поле України, ЄС і міжнародне, аналізує кримінально-правову протидію, стандарти управління інцидентами й технічні бази знань, механізми оцінювання відповідності та завершується кейсом визначення вимог для оператора критичної інфраструктури. Процеси реагування й криміналістики, що спираються на ці норми, розглянуто в розділі~\ref{chap:ch16}.

\section{Структура нормативного поля}\label{sec:ch02-structure}

Нормативне забезпечення протидії злоякісному коду доцільно подавати як ієрархію рівнів (рис.~\thechapter.1), що відрізняються за юридичною силою й предметом.

\emph{Міжнародний рівень} формують конвенції й угоди, насамперед Конвенція Ради Європи про кіберзлочинність (Будапештська конвенція)~\cite{budapest}, що гармонізує склади злочинів і механізми міжнародної співпраці, та її Другий додатковий протокол щодо транскордонного доступу до електронних доказів~\cite{ch02:budapest-ap}.

\emph{Наднаціональний рівень} для України визначається зобов'язаннями за Угодою про асоціацію з ЄС~\cite{ch02:assoc} і орієнтацією на право ЄС: Директиву NIS~2~\cite{nis2}, Акт про кіберстійкість~\cite{cra}, Акт про кібербезпеку~\cite{ch02:csa}, GDPR~\cite{gdpr}.

\emph{Національний законодавчий рівень} становлять закони України (про основні засади забезпечення кібербезпеки, про критичну інфраструктуру, про захист інформації в ІКС, про захист персональних даних), Кримінальний кодекс~\cite{kk-ukr} і підзаконні акти.

\emph{Рівень стандартів} містить міжнародні й національні стандарти (ISO/IEC~27000, NIST SP, ДСТУ), що набувають обов'язковості через посилання в нормативних актах або договорах.

\emph{Технічний рівень} формують бази знань і специфікації (MITRE~ATT\&CK, CVE, CVSS, STIX/TAXII, YARA, Sigma), які де-факто є спільною мовою галузі.

\begin{tikzfig}
\begin{tikzpicture}[font=\scriptsize,
  lv/.style={draw,rounded corners=2pt,minimum height=6.5mm,align=center},
  >={Stealth[length=1.6mm]}]
\node[lv,minimum width=38mm] (l1) at (0,0) {Міжнародний: Будапештська конвенція};
\node[lv,minimum width=50mm] (l2) at (0,-0.9) {Наднаціональний: право ЄС (NIS~2, CRA, GDPR)};
\node[lv,minimum width=62mm] (l3) at (0,-1.8) {Національний: закони України, КК, акти КМУ, НД~ТЗІ};
\node[lv,minimum width=74mm] (l4) at (0,-2.7) {Стандарти: ISO/IEC~27000, NIST SP, ДСТУ};
\node[lv,minimum width=86mm] (l5) at (0,-3.6) {Технічний: ATT\&CK, CVE/CVSS, STIX/TAXII, YARA, Sigma};
\node[rotate=90,anchor=south] at (4.7,-1.8) {\scriptsize деталізація $\downarrow$};
\end{tikzpicture}
\figcap{Рис.~\thechapter.1 --- Рівні нормативного забезпечення протидії злоякісному коду}
\end{tikzfig}

\section{Кримінально-правова протидія}\label{sec:ch02-criminal}

Створення, поширення й використання злоякісного коду є кримінальним правопорушенням. У Кримінальному кодексі України~\cite{kk-ukr} склади злочинів у сфері використання ЕОМ, систем та мереж електрозв'язку зосереджено в розділі~XVI. Ключовою для предмета є \term{стаття~361-1}~\cite{ch02:kk361} --- «Створення з метою використання, поширення або збуту шкідливих програмних чи технічних засобів, а також їх поширення або збут»: вона встановлює відповідальність саме за дії зі шкідливим кодом як таким, незалежно від настання наслідків. Стаття~361 карає несанкціоноване втручання в роботу інформаційних систем (типовий наслідок застосування ЗПК), статті~362--363-1 --- несанкціоновані дії з інформацією, її перехоплення й перешкоджання роботі мереж. Процесуальні аспекти (вилучення, фіксація й дослідження електронних доказів) регулює Кримінальний процесуальний кодекс~\cite{ch02:kpk}, а методичні вимоги до роботи з цифровими доказами --- стандарт ISO/IEC~27037~\cite{iso-27037} (розділ~\ref{chap:ch16}).

Важливою для фахівців є проблема \emph{правомірності дослідження} злоякісного коду. Створення, зберігання й аналіз зразків ЗПК є невід'ємною частиною роботи антивірусних лабораторій, дослідників безпеки, спеціалістів з реагування на інциденти та викладачів. Стаття~361-1 криміналізує дії, вчинені \emph{з метою} використання, поширення або збуту шкідливих засобів; отже, визначальним є намір (mens rea), а не сам факт володіння зразком чи його дослідження в захищеному середовищі з науковою або захисною метою. Проте межа тонка, тому дослідницька й освітня діяльність має спиратися на чіткі організаційні засади: ізольовані лабораторні середовища, контроль доступу до колекцій зразків, документування мети робіт, заборону поширення працездатних зразків поза контрольованим контуром, етичні норми відповідального розкриття вразливостей. У цьому посібнику виклад має суто захисну й аналітичну спрямованість: розглядаються принципи й методи \emph{протидії} та \emph{виявлення}, а не створення працездатного шкідливого коду. Такий підхід відповідає і міжнародній практиці, де законодавство про «хакерські інструменти» (dual-use) супроводжується застереженнями щодо законного застосування засобів тестування безпеки й дослідження.

Специфіка кримінально-правової протидії ЗПК --- \emph{транскордонність}. Зловмисник, інфраструктура й жертва часто перебувають у різних юрисдикціях, а докази --- у хмарних сервісах третіх країн. Будапештська конвенція~\cite{budapest} гармонізує склади злочинів і запроваджує механізми взаємної правової допомоги й збереження даних, а її Другий протокол~\cite{ch02:budapest-ap} спрощує транскордонний доступ до електронних доказів. Проте дієвість цих механізмів обмежена, якщо держава походження атаки не співпрацює, --- саме тому для державних APT (розділ~\ref{chap:ch01}) кримінально-правове стримування практично не працює, і наголос зміщується на технічний захист та атрибуцію (розділ~\ref{chap:ch16}).

\section{Законодавство України про кібербезпеку}\label{sec:ch02-ukraine}

Базовим актом є Закон України «Про основні засади забезпечення кібербезпеки України»~\cite{zu-cyber}, що визначає суб'єктів національної системи кібербезпеки, об'єкти кіберзахисту й повноваження. Для протидії ЗПК ключовими є обов'язки щодо кіберзахисту об'єктів критичної інфраструктури, повідомлення про кіберінциденти й взаємодії з CERT-UA~\cite{cert-ua}. Закон «Про критичну інфраструктуру»~\cite{zu-ki} встановлює категоризацію об'єктів і вимоги до їхнього захисту, деталізовані Загальними вимогами до кіберзахисту~\cite{kmu-518,ch02:kmu-943}, які прямо охоплюють захист від шкідливого коду: антивірусний захист, контроль застосунків, сегментацію, моніторинг, реагування на інциденти й відновлення.

Закон «Про захист інформації в інформаційно-комунікаційних системах»~\cite{zu-zi} і порядок створення комплексних систем захисту інформації (НД~ТЗІ~3.7-003-05~\cite{ndtzi-3-7-003}) визначають вимоги до захисту інформації з обмеженим доступом, зокрема до засобів захисту від несанкціонованого доступу й шкідливого коду; критерії оцінки захищеності встановлює НД~ТЗІ~2.5-004-99~\cite{ndtzi-2-5-004}, а термінологію --- НД~ТЗІ~1.1-003-99~\cite{ndtzi-1-1-003}. Важливо, що для систем, де обробляється інформація з обмеженим доступом, вимога застосування засобів захисту з підтвердженою відповідністю (позитивним висновком державної експертизи) є обов'язковою, що впливає на вибір антивірусних і EDR-рішень: для таких систем потрібні засоби, що пройшли відповідне оцінювання, тоді як для загальних систем вибір ширший і визначається ризик-орієнтованим підходом. Закон «Про електронні комунікації»~\cite{ch02:zu-ek} додатково покладає на постачальників обов'язки щодо цілісності й безпеки мереж, зокрема протидії зловмисному трафіку й повідомлення про інциденти. Типове положення про службу захисту інформації~\cite{ch02:ndtzi-1-4} визначає організаційні засади. Стратегічні пріоритети, зокрема розвиток спроможностей протидії, сформульовано у Стратегії кібербезпеки України~\cite{ukaz-447}. Досвід масштабних атак (розділ~\ref{chap:ch01}) і наближення до NIS~2 зумовили оновлення законодавства в напрямі ризик-орієнтованого підходу, розширення кола суб'єктів і чітких строків повідомлення про інциденти.

\section{Право Європейського Союзу}\label{sec:ch02-eu}

\subsection{Директива NIS~2}

Директива (ЄС)~2022/2555~\cite{nis2} встановлює високий спільний рівень кібербезпеки в ЄС і є орієнтиром для гармонізації українського законодавства. Для протидії ЗПК ключовими є: обов'язок уживати відповідних і пропорційних заходів управління ризиками (ст.~21), серед яких політики оброблення інцидентів, неперервність діяльності, безпека ланцюга постачання, кібергігієна, використання криптографії; багатоетапне повідомлення про значні інциденти (ст.~23); персональна відповідальність керівництва (ст.~20); суттєві адміністративні штрафи. Графік повідомлення задає систему обмежень на моменти подання звітів $t_1,t_2,t_3$ відносно моменту $t_0$, коли суб'єкту стало відомо про значний інцидент:

\begin{equation}
t_1 - t_0 \le 24~\text{год},\qquad t_2 - t_0 \le 72~\text{год},\qquad t_3 - t_2 \le 1~\text{міс.}
\label{eq:ch02-report}
\end{equation}

Обмеження~\eqref{eq:ch02-report} мають прямі технічні наслідки: система виявлення й реагування (розділи~\ref{chap:ch11}, \ref{chap:ch16}) повинна забезпечувати середній час виявлення (MTTD) і первинної класифікації інциденту з запасом менший за 24~год. Для інцидентів зі шкідливим кодом це означає потребу в автоматизованому виявленні й відпрацьованих процедурах тріажу.

\subsection{Акт про кіберстійкість та інші акти}

Регламент (ЄС)~2024/2847 --- Акт про кіберстійкість (CRA)~\cite{cra} --- встановлює вимоги кібербезпеки до продуктів із цифровими елементами протягом життєвого циклу: безпека за замовчуванням, відсутність відомих експлуатованих вразливостей на момент випуску, захищене оновлення, оброблення вразливостей, специфікація складу програмного забезпечення (SBOM, розділ~\ref{chap:ch15}), повідомлення про активно експлуатовані вразливості. CRA є інструментом зменшення поверхні атаки на рівні екосистеми продуктів: він зменшує кількість вразливостей, які зловмисники використовують для доставки й виконання ЗПК, і підвищує прозорість складу ПЗ, що критично для протидії атакам на ланцюг постачання (розділ~\ref{chap:ch15}). Зобов'язання щодо повідомлення про активно експлуатовані вразливості застосовуються раніше за основні вимоги, що створює механізм раннього попередження для всієї екосистеми. Для України, що гармонізує законодавство з правом ЄС, CRA задає орієнтир для вимог до постачальників цифрових продуктів, зокрема для критичної інфраструктури. Акт про кібербезпеку~\cite{ch02:csa} заснував європейську рамку сертифікації, а GDPR~\cite{gdpr} встановлює обов'язки щодо повідомлення про витоки персональних даних, які часто є наслідком дії ЗПК. Європейський центр компетенцій з кібербезпеки~\cite{ch02:eccc} координує дослідження й розвиток спроможностей.

\section{Стандарти управління та реагування}\label{sec:ch02-standards}

Загальну рамку управління безпекою задають стандарти серії ISO/IEC~27000, прийняті в Україні як ДСТУ EN ISO/IEC~27001~\cite{dstu-27001} і 27002~\cite{dstu-27002}; останній містить конкретні заходи щодо захисту від шкідливого коду, керування вразливостями й технічними заходами. Управління інцидентами регламентує ISO/IEC~27035~\cite{iso-27035}, настанови з інтернет-безпеки --- ISO/IEC~27032~\cite{ch02:iso27032}, управління ризиками --- ISO/IEC~27005~\cite{iso-27005} та ISO~31000~\cite{iso-31000}, неперервність --- ISO~22301~\cite{iso-22301}. Рамка NIST~CSF~2.0~\cite{nist-csf2} з функціями «Керувати», «Ідентифікувати», «Захищати», «Виявляти», «Реагувати», «Відновлювати» зручна для побудови профілів. Спеціалізованими є настанови NIST: SP~800-83 щодо запобігання інцидентам зі шкідливим кодом та їх обробки~\cite{nist-83} і SP~800-61 щодо обробки інцидентів~\cite{nist-61}. Пріоритетний набір практичних заходів дають CIS Controls~\cite{ch02:cis}, зокрема засоби захисту від шкідливого коду, керування вразливостями й контроль застосунків. Настанова NIST~SP~800-83~\cite{nist-83} заслуговує окремої уваги як профільний документ: вона структурує протидію за чотирма групами заходів --- політика (що дозволено й заборонено), підвищення обізнаності (навчання користувачів розпізнавати фішинг і соціальну інженерію), пом'якшення вразливостей (оновлення, конфігурація, принцип найменших привілеїв) і пом'якшення загроз (антивірусні засоби, контроль застосунків, системи виявлення вторгнень). Вона ж деталізує фази обробки інциденту зі шкідливим кодом --- підготовку, виявлення й аналіз, стримування, викорінення, відновлення та висновки, --- що узгоджується з ISO/IEC~27035 і становить основу розділу~\ref{chap:ch16}. Важливо, що стандарти розрізняють \emph{запобіжні} заходи (зменшують імовірність і масштаб зараження) і \emph{детективно-реактивні} (забезпечують виявлення й реагування за припущення, що частина атак подолає запобіжні рубежі), і вимагають обох груп. Відповідність основних документів функціям протидії подано в табл.~\thechapter.1.

\tabhead{Таблиця~\thechapter.1 --- Нормативні документи за функціями протидії злоякісному коду}

\begin{center}\footnotesize
\begin{tabularx}{\linewidth}{L{26mm} Y}
\toprule
Функція & Документи \\
\midrule
Управління ризиками & ISO/IEC~27005, ISO~31000, NIST~CSF~2.0 \\
Захист (заходи) & ДСТУ EN ISO/IEC~27002, CIS Controls, NIST~SP~800-83 \\
Виявлення & NIST~SP~800-137, MITRE~ATT\&CK/CAR \\
Реагування на інциденти & ISO/IEC~27035, NIST~SP~800-61, НД~ТЗІ \\
Криміналістика & ISO/IEC~27037, NIST~SP~800-86 \\
Неперервність і відновлення & ISO~22301 \\
Безпека продуктів і ланцюга & CRA, NIST~SP~800-218, 800-161 \\
\bottomrule
\end{tabularx}
\end{center}

\section{Технічні бази знань і специфікації}\label{sec:ch02-technical}

Окремий рівень --- технічні бази знань, що забезпечують узгоджений опис загроз і обмін інформацією. База MITRE~ATT\&CK~\cite{mitre-attack} каталогізує тактики й техніки зловмисників, стаючи спільною мовою для опису поведінки ЗПК (розділ~\ref{chap:ch04}); MITRE~D3FEND~\cite{mitre-defend} описує контрзаходи, CAR~\cite{mitre-car} --- аналітику виявлення. Вразливості каталогізує CVE~\cite{cve}, їхню серйозність оцінює CVSS~\cite{ch02:cvss}, патерни атак --- CAPEC~\cite{capec}. Для обміну розвідданими про загрози стандартизовано формат STIX і протокол TAXII~\cite{stix,ch02:taxii}, а платформи на кшталт MISP~\cite{ch02:misp} забезпечують практичний обмін індикаторами. Для виявлення застосовують правила YARA (ознаки у файлах і пам'яті)~\cite{yara} і Sigma (ознаки в журналах)~\cite{sigma}. Ці інструменти детально розглянуто в розділах~\ref{chap:ch06}, \ref{chap:ch11} і~\ref{chap:ch17}.

Принципова перевага відкритих баз знань і форматів --- \emph{сумісність і колективна оборона}: індикатори, виявлені однією організацією, швидко поширюються й використовуються іншими, що особливо важливо для протидії масштабним кампаніям (розділ~\ref{chap:ch01}). Водночас обмін має враховувати чутливість даних і довіру між учасниками, що регулюється моделями поширення (наприклад, світлофорний протокол TLP).

Корисною концептуальною рамкою для оцінювання цінності індикаторів є «піраміда болю»~\cite{pyramid}: що вищий рівень індикатора, тим дорожче зловмиснику його змінити й тим більшої шкоди завдає його блокування. У основі піраміди --- геш-значення файлів (їх зловмисник змінює тривіально), далі --- IP-адреси й доменні імена (змінюються легко), мережеві й хостові артефакти, засоби, і нарешті --- тактики, техніки й процедури (TTP), змінити які найважче, бо вони відображають спосіб мислення й робочі звички зловмисника. Ця рамка пояснює стратегічну перевагу поведінкового виявлення (розділ~\ref{chap:ch07}) над суто сигнатурним (розділ~\ref{chap:ch06}): виявлення за TTP, закодованими, зокрема, у правилах Sigma й базі ATT\&CK, витісняє зловмисника на рівень, де адаптація коштує йому найдорожче. Нормативно обмін розвідданими заохочується NIS~2 (обмін інформацією про загрози між суб'єктами) і реалізується через національні й галузеві майданчики за координації CERT-UA~\cite{cert-ua}.

\section{Оцінювання відповідності та обов'язки суб'єктів}\label{sec:ch02-compliance}

Для практичного аналізу стану відповідності множині вимог використовують кількісні показники. Нехай $\mathcal{R}=\{r_1,\dots,r_N\}$ --- множина застосовних вимог, $w_i>0$ --- вага вимоги (за критичністю), $c_i\in[0;1]$ --- ступінь її виконання. Тоді \term{зважений індекс розриву відповідності}

\begin{equation}
G = 1 - \frac{\sum_{i=1}^{N} w_i\, c_i}{\sum_{i=1}^{N} w_i},\qquad G\in[0;1].
\label{eq:ch02-gap}
\end{equation}

Індекс~\eqref{eq:ch02-gap} є адитивною згорткою й тому компенсаторний: повне виконання багатьох другорядних вимог може замаскувати невиконання критичної. Для обов'язкових вимог застосовують додаткову блокувальну умову $\min_{i:\,r_i\in\mathcal{R}_{\text{обов}}} c_i = 1$, тобто відповідність визнається лише за повного виконання кожної обов'язкової вимоги. Такий двоступеневий підхід узгоджується з логікою оцінювання ризиків (розділ~\ref{chap:ch04}) і дозволяє пріоритизувати усунення розривів за критерієм «зменшення $G$ на одиницю витрат».

Суб'єктами обов'язків є: оператори об'єктів критичної інфраструктури (кіберзахист, повідомлення, взаємодія з CERT-UA); власники ІКС з інформацією з обмеженим доступом (КСЗІ); постачальники цифрових продуктів (вимоги CRA); оператори персональних даних (захист від витоків). Одна організація може належати до кількох категорій одночасно.

Окремо слід наголосити на зв'язку відповідності й реальної захищеності. Формальне виконання вимог (наявність антивірусу, задокументована політика) не тотожне ефективній протидії: засіб може бути некоректно налаштований, сигнатури --- застарілі, моніторинг --- без реагування, резервні копії --- доступні зловмиснику з тієї самої мережі. Тому сучасні підходи доповнюють перевірку відповідності \emph{перевіркою ефективності}: моделюванням атак (red team, purple team), валідацією засобів контролю за реальними техніками ATT\&CK, навчаннями з реагування на кіберполігонах~\cite{ch02:polotai-ids}. Це узгоджується з переходом від «паперової» відповідності до ризик-орієнтованого й доказового підходу, закладеного в NIS~2 і NIST~CSF~2.0. Для протидії ЗПК така перевірка особливо важлива, оскільки зловмисники цілеспрямовано тестують свій код проти поширених засобів захисту (розділ~\ref{chap:ch08}), і формальна наявність засобу не гарантує виявлення конкретної загрози.

\section{Порівняння підходів України та ЄС}\label{sec:ch02-compare}

Порівняння ключових аспектів регулювання наведено в табл.~\thechapter.2. Воно показує, що національне законодавство рухається в напрямі, заданому правом ЄС, із збереженням специфіки (обов'язкова державна експертиза засобів, система ТЗІ).

\tabhead{Таблиця~\thechapter.2 --- Порівняння підходів України та ЄС до протидії злоякісному коду}

\begin{center}\footnotesize
\begin{tabularx}{\linewidth}{L{20mm} Y Y}
\toprule
Аспект & Україна & ЄС \\
\midrule
Рамковий акт & ЗУ «Про основні засади забезпечення кібербезпеки» & Директива NIS~2 \\
Криміналізація ЗПК & КК, ст.~361-1 та суміжні & національне право за Будапештською конвенцією \\
Повідомлення про інциденти & CERT-UA, строки за підзаконними актами & 24~год / 72~год / 1~міс. \\
Безпека продуктів & загальні вимоги & Акт про кіберстійкість (CRA) \\
Міжнародна співпраця & Будапештська конвенція, CERT-мережі & ENISA, мережа CSIRT, спільні операції \\
\bottomrule
\end{tabularx}
\end{center}

\section{Кейс: правові вимоги для оператора критичної інфраструктури}\label{sec:ch02-case}

Розглянемо оператора об'єкта критичної інфраструктури (наприклад, енергетичної компанії), що визначає застосовні вимоги протидії ЗПК. Процедура складається з кроків.

\begin{enumerate}
\item \emph{Ідентифікація ролей}: об'єкт критичної інфраструктури та об'єкт кіберзахисту (ЗУ «Про критичну інфраструктуру», «Про основні засади забезпечення кібербезпеки»); власник ІКС, можливо з інформацією з обмеженим доступом (ЗУ «Про захист інформації в ІКС»); оператор персональних даних.
\item \emph{Формування реєстру вимог}: обов'язкові заходи кіберзахисту~\cite{kmu-518,ch02:kmu-943} (антивірусний захист, контроль застосунків, сегментація, моніторинг, резервне копіювання, реагування), вимоги до повідомлення про інциденти, вимоги ТЗІ.
\item \emph{Відображення на технічні заходи}: вимога захисту від шкідливого коду --- на засоби захисту кінцевих точок і мережі (розділи~\ref{chap:ch10}, \ref{chap:ch11}); вимога виявлення --- на моніторинг і EDR; вимога неперервності --- на ізольовані резервні копії й план відновлення (розділ~\ref{chap:ch16}); вимога повідомлення --- на процедуру з MTTD у межах~\eqref{eq:ch02-report}.
\item \emph{Оцінювання розриву}: розрахунок індексу~\eqref{eq:ch02-gap} з перевіркою блокувальної умови для обов'язкових вимог.
\item \emph{Планування}: пріоритизація заходів за ризиком і вартістю (розділ~\ref{chap:ch04}).
\end{enumerate}

Числовий приклад. Реєстр містить 30 вимог: 10 обов'язкових із вагою~3, 14 стандартів із вагою~2, 6 рекомендацій із вагою~1 (сума ваг~64). Виконано: обов'язкових --- 8 повністю, 2 частково (по~0,5); стандартів --- 10 повністю, 4 частково; рекомендацій --- 3 повністю. Зважена сума: $3(8{+}1){+}2(10{+}2){+}1\cdot3 = 27{+}24{+}3 = 54$, звідки $G = 1 - 54/64 \approx 0{,}156$. Попри помірний індекс, блокувальна умова не виконана через дві частково реалізовані обов'язкові вимоги (приміром, відсутність ізольованих резервних копій і неповне покриття моніторингом), тож саме вони першочергові, навіть якщо усунення інших розривів дало б більше зменшення $G$. Приклад ілюструє, що для об'єктів критичної інфраструктури формальний індекс доповнюється безумовним виконанням критичних вимог, спрямованих на найруйнівніші сценарії (розділ~\ref{chap:ch01}). Інтеграцію засобів виявлення й реагування та перевірку відповідних процедур на кіберполігонах розглянуто у вітчизняних дослідженнях~\cite{ch02:polotai-ids}, а підходи до класифікації загроз --- у~\cite{ch02:shevchenko-class}.

\section*{Висновки до розділу}
\addcontentsline{toc}{section}{Висновки до розділу}

Протидія злоякісному коду діє в багаторівневому нормативному полі: міжнародні конвенції гармонізують кримінальну відповідальність і співпрацю, право ЄС і національне законодавство встановлюють обов'язки організацій, стандарти конкретизують процеси, а технічні бази знань забезпечують спільну мову. Кримінально-правова протидія (в Україні --- насамперед стаття~361-1 КК, що карає дії зі шкідливим кодом як такі) ускладнена транскордонністю й неефективна проти державних порушників, що зміщує наголос на технічний захист і атрибуцію. Законодавство України про кібербезпеку й критичну інфраструктуру прямо зобов'язує до захисту від шкідливого коду, виявлення, реагування й відновлення, а наближення до NIS~2 запроваджує ризик-орієнтований підхід і строки повідомлення про інциденти~\eqref{eq:ch02-report}, що визначають вимоги до MTTD. Стандарти ISO/IEC~27000, NIST~SP~800-83/61, CIS~Controls задають процеси й заходи, а відкриті бази знань (ATT\&CK, CVE/CVSS, STIX/TAXII, YARA, Sigma) --- технічну основу колективної оборони. Для практики оцінювання корисний двоступеневий підхід: блокувальні обов'язкові вимоги плюс зважений індекс розриву~\eqref{eq:ch02-gap}, причому для критичної інфраструктури пріоритет надається безумовному виконанню вимог, спрямованих на найруйнівніші сценарії.

\selfcheck
\begin{enumerate}
\item Назвіть рівні нормативного забезпечення протидії злоякісному коду (рис.~\thechapter.1) і наведіть документи кожного рівня.
\item Яку відповідальність встановлює стаття~361-1 КК України? Чим вона відрізняється від статті~361 за об'єктом?
\item У чому полягає транскордонна специфіка кримінально-правової протидії ЗПК? Які механізми пропонує Будапештська конвенція і чому вони не працюють проти державних APT?
\item Які обов'язки щодо протидії шкідливому коду встановлює законодавство України про критичну інфраструктуру?
\item Які заходи управління ризиками вимагає стаття~21 Директиви NIS~2? Як строки повідомлення~\eqref{eq:ch02-report} впливають на вимоги до MTTD?
\item Що встановлює Акт про кіберстійкість і як він зменшує поверхню атаки на рівні екосистеми продуктів?
\item Зіставте нормативні документи з функціями протидії (табл.~\thechapter.1). Які документи стосуються криміналістики й реагування?
\item Яку роль відіграють технічні бази знань (ATT\&CK, CVE/CVSS, STIX/TAXII, YARA, Sigma)? У чому перевага відкритих форматів для колективної оборони?
\item Обчисліть індекс розриву~\eqref{eq:ch02-gap} для прикладу підрозділу~\ref{sec:ch02-case}, якщо обидві частково виконані обов'язкові вимоги буде виконано повністю. Чи виконується тепер блокувальна умова?
\item Хто є суб'єктами обов'язків щодо протидії ЗПК? Чому одна організація може належати до кількох категорій?
\item Порівняйте підходи України та ЄС (табл.~\thechapter.2). Які напрями гармонізації ви вважаєте першочерговими?
\item Чому для об'єктів критичної інфраструктури формальний індекс відповідності доповнюється безумовним виконанням критичних вимог?
\item Як стандарт ISO/IEC~27037 і НД~ТЗІ пов'язані з кримінальним процесом щодо інцидентів зі шкідливим кодом?
\item Дослідницьке: розробіть реєстр вимог протидії ЗПК для організації обраного сектору, зіставте вимоги з технічними заходами й обчисліть індекс розриву з перевіркою блокувальних умов.
\item Дослідницьке: проаналізуйте застосовність механізмів Другого протоколу до Будапештської конвенції для отримання електронних доказів у справах про поширення ЗПК і визначте обмеження за наявної міжнародної практики.
\end{enumerate}
